\begin{abstract}
Robot learning is splitting into two bets: policies that bake competence into \emph{frozen
weights} (vision-language-action, or VLA, models), and agents that write and refine their own
\emph{executable skills} as code. This survey organises the field around that axis of
\textbf{weights versus skills}. Its central analytical contribution is a deep-dive that arranges
code-as-policy methods by their \emph{degree of self-improvement}, from zero-shot program
synthesis, through closed-loop self-repair and persistent skill memory, to the sparsely
populated cell in which execution feedback, skill memory, and evolutionary search combine into
one open-ended loop; only a few very recent systems (for example ASPIRE, ENPIRE, and RoboClaw)
occupy that cell. We map the complementary ``skills'' pole, from unsupervised
reinforcement-learning skill discovery to large-language-model skill libraries, and show that
the word ``skill'' is used in at least five distinct senses, of which only the code sense
self-improves without gradient updates. We then connect the taxonomy to the emerging
\emph{skill economy}: commercial robot-skill marketplaces now distribute one-tap skills across
robots but ship only static playback, which surfaces open problems of adaptation,
cross-embodiment portability, provenance, safety verification, composition, and
standardisation. This is a deliberately \emph{focused} survey. Rather than cataloguing the
field exhaustively, it examines 77 representative systems across six technique families through
one taxonomy and a set of contrast tables, and it supplies operational definitions of the
self-improvement mechanisms together with a statement of what each family cannot do.
\end{abstract}
